\section{基于转移的依存分析（Transition-based Parsing）}

\subsection{依存分析的方法概览}

依存分析有多种算法范式：

\begin{enumerate}
    \item \textbf{动态规划方法}（Dynamic Programming）
    \item \textbf{基于图的方法}（Graph-based / Maximum Spanning Tree）
    \item \textbf{约束满足方法}（Constraint Satisfaction）
    \item \textbf{基于转移的方法}（Transition-based / Shift-Reduce）——本讲重点
\end{enumerate}

基于转移的方法在实践中最为流行，因为它具有\textbf{线性时间复杂度}，远优于上下文无关文法分析的立方时间复杂度。

\subsection{Arc-Standard 转移系统}

基于转移的依存分析使用两个核心数据结构和三种操作：

\begin{importantbox}{转移系统的组成}
\textbf{数据结构}：
\begin{itemize}
    \item \textbf{栈}（Stack）：存放正在处理的词，初始只包含 ROOT。\textbf{栈顶在右侧}
    \item \textbf{缓冲区}（Buffer）：存放尚未处理的词，初始包含句子的所有词。\textbf{缓冲区顶在左侧}
\end{itemize}

\textbf{三种操作}：
\begin{enumerate}
    \item \textbf{SHIFT}：将缓冲区顶部的词移到栈顶
    \item \textbf{LEFT-ARC}：栈顶第二个词成为栈顶词的依存词（左弧），弹出第二个词
    \item \textbf{RIGHT-ARC}：栈顶词成为栈顶第二个词的依存词（右弧），弹出栈顶词
\end{enumerate}

\textbf{终止条件}：缓冲区为空，栈中只剩 ROOT。
\end{importantbox}

\subsection{实例演示：分析 ``I ate fish''}

\begin{figure}[H]
\centering
\includegraphics[width=0.95\textwidth]{slides-images/slide-028.jpg}
\caption{转移系统的操作定义与初始/终止条件\protect\footnotemark}
\end{figure}
\footnotetext{Slides 第28页。}

以句子 ``I ate fish'' 为例，分析过程如下：

\begin{table}[H]
\centering
\small
\begin{tabular}{clll}
\toprule
\textbf{步骤} & \textbf{栈} & \textbf{缓冲区} & \textbf{操作} \\
\midrule
0 & [ROOT] & [I, ate, fish] & 初始状态 \\
1 & [ROOT, I] & [ate, fish] & SHIFT \\
2 & [ROOT, I, ate] & [fish] & SHIFT \\
3 & [ROOT, ate] & [fish] & LEFT-ARC（I $\leftarrow$ ate） \\
4 & [ROOT, ate, fish] & [] & SHIFT \\
5 & [ROOT, ate] & [] & RIGHT-ARC（ate $\rightarrow$ fish） \\
6 & [ROOT] & [] & RIGHT-ARC（ROOT $\rightarrow$ ate） \\
\bottomrule
\end{tabular}
\caption{``I ate fish'' 的转移操作序列}
\end{table}

\begin{figure}[H]
\centering
\includegraphics[width=0.95\textwidth]{slides-images/slide-029.jpg}
\caption{``I ate fish'' 的逐步转移过程\protect\footnotemark}
\end{figure}
\footnotetext{Slides 第29页。}

最终构建的依存关系集合为：
\begin{itemize}
    \item ate $\rightarrow$ I（主语）
    \item ate $\rightarrow$ fish（宾语）
    \item ROOT $\rightarrow$ ate（句子根节点）
\end{itemize}

\begin{warningbox}{操作选择决定了分析结果}
不同的操作序列会生成不同的依存结构。如果在第 3 步选择 RIGHT-ARC 而非 LEFT-ARC，就会让 ``I'' 成为中心词而 ``ate'' 成为依存词——这显然是错误的。因此，\textbf{如何在每一步选择正确的操作}是整个系统的核心问题。
\end{warningbox}

\subsection{Nivre 与贪心决策}

\begin{figure}[H]
\centering
\includegraphics[width=0.95\textwidth]{slides-images/slide-030.jpg}
\caption{Joakim Nivre 与 MaltParser\protect\footnotemark}
\end{figure}
\footnotetext{Slides 第30页。}

瑞典 NLP 研究者 \textbf{Joakim Nivre}（2003--2005）提出了用\textbf{机器学习}来预测每一步应采取的操作。核心洞察：

\begin{importantbox}{贪心转移分析的关键发现}
\begin{enumerate}
    \item 每步操作都可以视为一个\textbf{分类问题}：给定当前状态（栈、缓冲区内容），预测应执行 SHIFT、LEFT-ARC 还是 RIGHT-ARC
    \item 尽管使用\textbf{贪心策略}（每步只选最优操作，不做搜索），分析精度仍然很高
    \item 整个分析过程是\textbf{线性时间} $O(n)$：每个词最多被 SHIFT 一次、参与一次 ARC 操作
\end{enumerate}
\end{importantbox}

\subsection{传统特征工程方法}

Nivre 最初（2005 年）使用的是传统的\textbf{符号化特征}（indicator features）配合线性分类器（如逻辑回归或 SVM）：

\begin{figure}[H]
\centering
\includegraphics[width=0.95\textwidth]{slides-images/slide-033.jpg}
\caption{传统方法的特征工程：使用大量手工设计的指示特征\protect\footnotemark}
\end{figure}
\footnotetext{Slides 第33页。}

典型特征包括：
\begin{itemize}
    \item 栈顶词是 ``good'' 且词性是形容词
    \item 栈顶词是 ``good'' 且栈中第二个词是 ``had''
    \item 缓冲区顶部词的词性是名词
\end{itemize}

这些特征的组合可以达到数百万甚至更多。

\begin{warningbox}{传统特征方法的三大问题}
\begin{enumerate}
    \item \textbf{稀疏性}（Sparsity）：大多数特征组合在训练数据中极少出现，可能仅出现 10 次/百万句
    \item \textbf{不完备性}（Incompleteness）：有些词和组合在训练数据中从未出现
    \item \textbf{计算开销}：运行时分析中，\textbf{95\% 以上的时间}花在计算特征上，而非机器学习决策本身
\end{enumerate}
\end{warningbox}

\subsection{本章小结}

基于转移的依存分析通过栈和缓冲区模拟分析过程，使用 SHIFT、LEFT-ARC、RIGHT-ARC 三种操作逐步构建依存树。Nivre 提出的贪心策略实现了线性时间复杂度。传统方法依赖大量手工特征，存在稀疏性、不完备性和高计算开销等问题。

%% ========== 神经依存分析 ==========

